% sections/generation_pipeline.tex — Generation Pipeline
\section{Generation Pipeline}
\label{sec:generation_pipeline}

\system{}'s generation pipeline transforms a declarative YAML
configuration into a complete, internally consistent enterprise
dataset---producing the \emph{structural knowledge layer}
($\mathcal{K}_S$) of the reference knowledge graph.

% ─────────────────────────────────────────────────────────────────────
\subsection{Orchestration Workflow}
\label{sec:gen_orchestration}

The \texttt{GenerationOrchestrator} executes the following stages:

\begin{enumerate}
  \item \textbf{Configuration loading and validation.}  Parse YAML,
    apply industry preset defaults, validate invariants.

  \item \textbf{Chart of Accounts initialization.}  Generate a GL
    account structure at the configured complexity level (small:
    $\sim$100, medium: $\sim$400, large: $\sim$2{,}500 accounts).

  \item \textbf{Master data generation.}  Produce vendors, customers,
    materials, fixed assets, employees, and cost centers with
    cross-referenced relationships.

  \item \textbf{Opening balance generation.}  Establish period-start
    balances satisfying the accounting equation $A = L + E$.

  \item \textbf{Document-flow generation.}  Execute P2P, O2C, S2C, HR,
    manufacturing, tax, treasury, ESG, and project accounting generators.

  \item \textbf{Period-close processing.}  Run accruals, depreciation,
    FX translation, intercompany elimination, and year-end closing.
    Generate trial balances and financial statements.

  \item \textbf{Standards generation.}  Produce accounting standards
    artifacts (ASC~606 contracts, ASC~842 leases, fair value measurements,
    impairment tests) and audit standards artifacts (ISA mappings,
    confirmations, audit opinions, SOX assessments).

  \item \textbf{Anomaly injection.}  Apply configured anomalies with
    ground-truth labels across all five knowledge dimensions.

  \item \textbf{Data quality variation.}  Optionally inject missing
    values, format variations, typos, and encoding issues.

  \item \textbf{Evaluation and graph export.}  Run coherence validators,
    construct the reference knowledge graph, and export to configured
    formats.

  \item \textbf{Output.}  Stream results to configured sinks (CSV, JSON,
    Parquet).
\end{enumerate}

% ─────────────────────────────────────────────────────────────────────
\subsection{Journal Entry Generation}
\label{sec:gen_journal}

Journal entries are the atomic unit of accounting data.

\begin{definition}[Balanced Entry]
A journal entry $j$ with line items $\{l_1, \ldots, l_n\}$ is
\emph{balanced} iff:
\begin{equation}
\sum_{i : l_i.\text{debit}} l_i.\text{amount}
  = \sum_{i : \neg l_i.\text{debit}} l_i.\text{amount}.
\label{eq:balanced}
\end{equation}
\end{definition}

\system{} enforces \Cref{eq:balanced} at construction time: the
\texttt{JournalEntry} constructor returns an error if the debit-credit
difference exceeds a configurable tolerance (default: \$0.01).  This
invariant is part of the \emph{normative knowledge layer} and is
maintained across all generators.

Line-item amounts are drawn from the configured mixture distribution
(\Cref{eq:lognormal_mixture}).  The number of line items per entry
follows a correlated draw from the copula model, ensuring that
higher-amount entries tend to have more line items and require higher
approval levels.

% ─────────────────────────────────────────────────────────────────────
\subsection{Document Flow Generation}
\label{sec:gen_docflow}

Enterprise transactions span multiple documents with strict sequencing
and referential integrity.

\paragraph{Procure-to-Pay (P2P).}
\[
\text{Purchase Order} \to \text{Goods Receipt} \to \text{Vendor Invoice}
\to \text{Payment}
\]
Each stage creates the corresponding document record, updates the
document-chain manager with forward/backward references, and posts
journal entries to the appropriate GL accounts.

\paragraph{Three-Way Matching.}
The three-way matcher validates consistency across PO, GR, and invoice
with quantity tolerance ($\epsilon_q = 2\%$), price tolerance
($\epsilon_p = 5\%$), and over-delivery allowance ($\delta_{\max} = 10\%$).

\paragraph{Order-to-Cash (O2C).}
\[
\text{Sales Order} \to \text{Delivery} \to \text{Customer Invoice}
\to \text{Customer Receipt}
\]
with revenue recognition, AR postings, and cash receipts.  Partial
deliveries and split invoicing are supported.

\paragraph{Source-to-Contract (S2C).}
Full sourcing lifecycle: sourcing projects, supplier qualifications,
RFx events, supplier bids, bid evaluations, procurement contracts,
catalog items, and supplier scorecards.

\paragraph{HR/Payroll.}
Payroll runs with line items (base pay, overtime, deductions, taxes),
time entries, expense reports, and benefit enrollment (health, dental,
vision, retirement) with employer/employee contributions.

\paragraph{Manufacturing.}
Production orders, routing operations, quality inspections with
characteristics, cycle counts, multi-level BOMs with phantom assembly
support, and inventory movement tracking (receipt, issue, transfer,
return, scrap, adjustment).

% ─────────────────────────────────────────────────────────────────────
\subsection{Tax Accounting}
\label{sec:gen_tax}

\system{} generates comprehensive tax accounting data spanning multiple
jurisdictions:

\begin{itemize}
  \item \textbf{Tax jurisdictions and codes:} Multi-level jurisdiction
    hierarchies (federal, state, local) with associated tax codes,
    rates, and effective dates.
  \item \textbf{Tax lines and returns:} Per-transaction tax calculations,
    periodic tax return preparation with aggregated liabilities.
  \item \textbf{Tax provisions:} Current and deferred tax calculations
    under ASC~740/IAS~12, including temporary differences, deferred tax
    rollforward, and rate reconciliation.
  \item \textbf{Withholding tax:} Source-based withholding records with
    jurisdiction-specific rates and treaty provisions.
  \item \textbf{Uncertain tax positions:} UTP assessment with probability
    weighting and disclosure requirements.
  \item \textbf{Transfer pricing:} Intercompany pricing methods
    (comparable uncontrolled, resale, cost-plus, TNMM, profit split)
    with arm's-length validation.
\end{itemize}

% ─────────────────────────────────────────────────────────────────────
\subsection{Treasury Management}
\label{sec:gen_treasury}

The treasury module generates:

\begin{itemize}
  \item \textbf{Cash positioning:} Daily cash positions across accounts
    and entities with automated sweeping via cash pools.
  \item \textbf{Cash forecasting:} Rolling forecasts with multiple
    time horizons and confidence intervals.
  \item \textbf{Hedging instruments:} Forward contracts, options, and
    swaps with hedge relationship documentation and effectiveness
    testing under ASC~815/IFRS~9.
  \item \textbf{Debt instruments:} Term loans, revolving facilities, and
    bonds with amortization schedules and covenant tracking.
  \item \textbf{Debt covenants:} Financial covenant calculations with
    compliance monitoring and breach detection.
\end{itemize}

% ─────────────────────────────────────────────────────────────────────
\subsection{ESG and Sustainability}
\label{sec:gen_esg}

The ESG module generates sustainability data aligned with major
reporting frameworks:

\begin{itemize}
  \item \textbf{Environmental:} GHG emissions (Scope~1/2/3), energy
    consumption by source, water usage, waste records by type and
    disposition, and climate scenario analysis.
  \item \textbf{Social:} Workforce diversity metrics, pay equity
    analysis, safety incidents and metrics, materiality assessments.
  \item \textbf{Governance:} Board composition metrics, supplier ESG
    assessments, ESG disclosure documents.
\end{itemize}

% ─────────────────────────────────────────────────────────────────────
\subsection{Project Accounting}
\label{sec:gen_project}

Project-based organizations require WBS-structured cost tracking:

\begin{itemize}
  \item \textbf{Projects:} Multi-level project structures with budget
    allocation and resource planning.
  \item \textbf{Cost lines:} Labor, material, and overhead cost
    postings with project/WBS element attribution.
  \item \textbf{Revenue recognition:} Percentage-of-completion and
    milestone-based revenue for long-term contracts.
  \item \textbf{Earned value:} BCWS, BCWP, ACWP calculations with
    derived metrics (CPI, SPI, EAC, ETC).
  \item \textbf{Change orders:} Scope and budget modifications with
    approval tracking.
  \item \textbf{Milestones:} Deliverable tracking with planned vs.\
    actual completion dates.
\end{itemize}

% ─────────────────────────────────────────────────────────────────────
\subsection{Subledger and Period-Close Processing}
\label{sec:gen_subledger}

\paragraph{Subledger generation.}
Four subledger generators maintain detailed subsidiary records:
AR (customer invoices, receipts, aging, credit memos),
AP (vendor invoices, payments, aging, prepayments),
FA (capitalization, depreciation, disposals, impairment),
and Inventory (receipts, issues, valuation methods).
Each reconciles to its control account in the GL.

\paragraph{Period-close engine.}
Executes: (i)~accrual entries, (ii)~depreciation postings,
(iii)~FX translation with CTA postings, (iv)~intercompany
matching and elimination, (v)~year-end income-to-equity
reclassification, and (vi)~financial statement generation
(balance sheet, income statement, cash flow, changes in equity).

% ─────────────────────────────────────────────────────────────────────
\subsection{Intercompany Processing}
\label{sec:gen_intercompany}

Multi-entity datasets require IC transaction matching and elimination
for consolidated reporting.  The IC generator creates bilateral
transaction pairs.  The matching engine identifies matched pairs within
configurable tolerances.  The elimination generator produces
consolidation entries that net IC receivables/payables, IC
revenue/expense, and unrealized profit in inventory.

% ─────────────────────────────────────────────────────────────────────
\subsection{Accounting and Audit Standards}
\label{sec:gen_standards}

The normative knowledge layer ($\mathcal{K}_N$) is generated through:

\paragraph{Accounting standards.}
Revenue recognition (ASC~606/IFRS~15) with customer contracts and
performance obligations.  Leases (ASC~842/IFRS~16) with ROU assets and
lease liabilities.  Fair value measurements (ASC~820/IFRS~13) with
Level~1/2/3 hierarchy.  Impairment testing (ASC~360/IAS~36).  ECL models
and provision matrices.  Business combinations with purchase price
allocation.  Pension obligations and stock compensation.

\paragraph{Audit standards.}
ISA compliance (34 standards) with procedure mapping.  PCAOB standards
(19+) with cross-reference.  Analytical procedures (ISA~520),
confirmations (ISA~505), audit opinions (ISA~700/705/706/701), key
audit matters.  Group audit (ISA~600) with component auditor model.

\paragraph{Internal controls.}
COSO~2013 framework with 5 components, 17 principles, and maturity
levels.  12 transaction-level controls (C001--C060) and 6 entity-level
controls (C070--C081).  Segregation of duties enforcement.
SOX~302/404 assessments with deficiency classification and material
weakness identification.

\paragraph{Regulatory exports.}
FEC (French \emph{Fichier des \'{E}critures Comptables}) and GoBD
(German) audit-ready exports for jurisdiction-specific tax authority
requirements.
